\begin{proposition}[The torsion-free cotangent assertion for arbitrary singular supports in a torus fibration]\hypertarget{tag-01AX}{}\label{cdp:st:A.4.8}
\StatementRecord{cdp:st:A.4.8}{proved}{01AX}{\SeriesRef{cdp:st:D.1.9}, \SeriesRef{cdp:st:D.1.8}, \SeriesRef{cdp:eq:A.22}, \SeriesRef{cdp:st:A.4.7}}

Let $g:Y\to\PP^1$ be a proper surjective holomorphic map with connected fibres from a smooth compact complex threefold with $b_1(Y)=b_2(Y)=0$, and suppose its general fibre is a complex two-torus. Let $L_r=\exp(2\pi\mathrm i r\beta)\in\Pic^0(Y)$ be the transgression family constructed in \ref{cdp:st:D.1.9}. For every real irrational $r$ and every reduced fibre support $S=(Y_c)_{\rm red}$,
\[
 H^1(S,\mathbb A_{r,S})=0,
 \qquad
 H^0(S,\widetilde\Omega_S^1\otimes L_r|_S)=0,
 \tag{A.21}\hypertarget{tag-004D}{}\label{cdp:eq:A.21}
\]
where $\mathbb A_{r,S}$ is the unitary local system whose associated holomorphic line is $L_r|_S$. In particular this proves Equation~(11) of CDP20 for all singular reduced fibre supports in the torus case, with no assumption on their singularities or their scheme-theoretic multiplicities.
\end{proposition}


\begin{proof}

\begin{enumerate}[label={$\langle1\rangle$\arabic*.},leftmargin=*]
\item We specify all comparison maps used from \ref{cdp:st:D.1.9}. On a sufficiently small disc $\Delta$ about $c$, put $U=g^{-1}(\Delta)$ and $U^*=U\setminus S$. The integral section $\gamma$ of $R^1g_*\Z$ has a lift $\gamma_U\in H^1(U,\Z)$ by Leray over the contractible disc. Its image in $H^1(U,\OO_U)$ equals $\beta|_U$, by the coherent Leray sequence. Thus the restriction of the holomorphic line $L_r$ to $U$ is the line associated to the flat character
\[
 \pi_1(U)\longrightarrow S^1,\qquad
 a\longmapsto\exp\bigl(2\pi\mathrm i r\langle\gamma_U,a\rangle\bigr).
 \tag{A.22}\hypertarget{tag-004E}{}\label{cdp:eq:A.22}
\]
Denote this local system by $\mathbb A_U$. \ref{cdp:st:D.1.8}--\ref{cdp:st:D.1.9} prove that $\gamma_U$ restricts nontrivially to every nearby smooth fibre. Irrationality of $r$ therefore makes \eqref{cdp:eq:A.22} nontrivial on such a fibre $F$.

\item A nontrivial rank-one local system on the real torus underlying $F$ has zero cohomology in every degree. Explicitly, use the standard cellular cochain complex of the product of four circles with coefficients in the character. It is the Koszul complex for the four scalars $\chi(e_j)-1$. At least one scalar is nonzero; contraction in that factor, divided by this scalar, is a contracting homotopy. Therefore $H^q(F,\mathbb A_U|_F)=0$ for every $q$. The map $U^*\to\Delta^*$ is a differentiable fibre bundle. Its local-system Leray spectral sequence has zero fibre-cohomology sheaves, hence
$H^1(U^*,\mathbb A_U|_{U^*})=0$.

\item Inclusion induces a surjection $\pi_1(U^*)\to\pi_1(U)$: a loop in the smooth real six-manifold $U$ can be perturbed off the divisor, whose real codimension is at least two, relative to a base point in $U^*$. The induced map on first cohomology with these local coefficients is injective. For an explicit verification, identify first cohomology with crossed homomorphisms modulo coboundaries. If a crossed homomorphism $a:G\to M$ becomes $a(q(h))=(\chi(q(h))-1)v$ after pullback by a surjection $q:H\to G$, surjectivity gives $a(g)=(\chi(g)-1)v$ for every $g\in G$. Thus its original class was already zero. Applied to the displayed fundamental-group map, this proves $H^1(U,\mathbb A_U)=0$.

\item The argument applies after every sufficiently small shrinking of $\Delta$, with the restrictions of the same local system. Taking the direct limit and using proper base change identifies the resulting stalk with $H^1(S,\mathbb A_U|_S)$; scheme-theoretic nilpotents do not affect this topological local system. This proves the first vanishing in \eqref{cdp:eq:A.21}. The second is exactly \ref{cdp:st:A.4.7} applied to the restriction of \eqref{cdp:eq:A.22}. All singular values are treated by the same global parameter $r$, so the assertion is simultaneous.

\end{enumerate}
\end{proof}